\section{研究计划进度安排及预期目标}

\subsection{进度安排}
本课题的研究周期计划为6个月，具体进度安排如下：

计划进度将按“文献调研与开题—数据与底座建设—核心算法研发—系统集成评测—论文撰写答辩”的顺序推进。2025年11月至12月阶段主要完成RAG、知识图谱与大模型推理（CoT/GoT）等方向的文献调研，明确研究切入点与总体技术路线，并完成开题报告与答辩，同时搭建实验环境并配置必要的GPU与索引组件。2025年12月至2026年1月阶段聚焦数据与基础架构：收集并清洗多源异构数据（Wikipedia、HotpotQA及垂直领域PDF/表格），构建混合索引系统并完成Baseline RAG流水线，为后续算法对比提供统一基准。2026年2月至3月阶段开展核心算法设计与实现，重点研发自适应检索控制与图结构推理机制，并通过小规模实验持续验证与迭代。2026年4月至5月阶段完成系统集成与全面评测，开发可视化演示界面，在多跳问答基准上进行对比实验与消融分析，评估各模块的贡献与稳定性。2026年5月至6月阶段整理实验结果与图表，完成论文撰写、修改定稿与答辩准备。

\begin{table}[ht]
\centering
\caption{研究进度安排（建议版，可按实际起止月份调整）}
\label{tab:proposal_schedule}
\small
\setlength{\tabcolsep}{0.6em}
\begin{tabularx}{\linewidth}{|c|c|X|X|}
\hline
\textbf{阶段} & \textbf{时间} & \textbf{主要任务} & \textbf{阶段产出} \\
\hline
阶段1 & 2025.11--2025.12 & 文献调研与问题定义；确定研究路线与评测口径；完成开题材料与环境搭建 & 开题报告；文献综述初稿；实验环境与代码框架 \\
\hline
阶段2 & 2025.12--2026.01 & 数据收集与清洗分块；混合索引原型；Baseline RAG全链路打通 & 数据与索引底座；Baseline指标；可复现实验脚本 \\
\hline
阶段3 & 2026.02--2026.03 & 自适应检索与路由策略；重排序与查询扩展；图结构推理原型与消融实验 & 核心算法模块；消融与对比结果；阶段性报告 \\
\hline
阶段4 & 2026.04--2026.05 & 系统集成与性能优化；多跳问答/异构融合实验；可解释与验证机制完善 & 集成原型系统；评测报告与可视化；可解释输出样例 \\
\hline
阶段5 & 2026.05--2026.06 & 论文撰写、修改定稿；查重与材料整理；答辩准备与演示完善 & 毕业论文定稿；答辩PPT；可演示系统与复现材料 \\
\hline
\end{tabularx}
\end{table}

\subsection{预期目标}
通过本课题的研究，预期达到以下目标：

在技术层面，预期构建一套面向多源异构知识的RAG框架，实现文本、图谱与表格等数据的高效索引与联合检索，并支持在不同查询类型下的自适应检索与证据融合；在复杂推理方面，争取在HotpotQA等多跳问答数据集上相对Baseline取得可量化的提升，并形成可复现的评测与消融结论；在可信性与可解释性方面，系统应能输出可追溯的证据链与推理路径，并对关键事实提供验证结果，从而降低无依据回答的风险。

在学术与工程层面，预期完成本科毕业论文并通过答辩，同时交付一个具备Web演示能力的原型系统，支持用户上传或接入异构文档并进行复杂问答交互；并在整理核心算法与实验脚本后形成可复现的代码资产，必要时进一步抽象为可复用组件以便开源共享；在条件允许的情况下，争取将研究成果整理为论文进行投稿。

\textbf{交付成果：}为保证进度计划可验收、可追踪，本文将交付成果拆分为可检查的工程与文档资产，覆盖“数据—索引—算法—评测—展示—论文”全链路：
\begin{enumerate}
    \item \textbf{数据与证据资产：}多源异构数据样例集（文本/表格/图谱）与清洗分块脚本；证据定位方案与基础实体对齐配置（如别名归一/ID映射的原型规则）；典型难例清单与失败案例分析记录。
    \item \textbf{索引与检索底座：}向量/倒排/图索引构建与更新脚本；检索控制器与路由模块原型；检索日志与成本统计（调用次数、Token、延迟）。
    \item \textbf{推理与验证模块：}证据链/推理链构造模块原型（含可回放推理路径）；冲突检测与拒答策略；验证与自检模块及其触发规则（可配置）。
    \item \textbf{评测与复现材料：}统一评测脚本与结果汇总工具，支持对比与消融；指标报告（效果/成本/支撑率）；复现指南与环境依赖说明。
    \item \textbf{演示系统与论文材料：}可交互的Web演示界面（答案+证据+推理路径展示）；毕业论文定稿与答辩PPT；必要的归档材料。
\end{enumerate}
